% P06 | 3 Part A: Brownian motion and random walks | Brownian motion and the scaled random walk
% guide: dev_notes/paragraph_guides/P06_brownian_motion_and_the_scaled_random_walk.txt
A process $W=\{W_t\}_{t\ge0}$ is a Brownian motion on $(\Omega,\mathcal F,\{\mathcal F_t\},\mathbb P)$ if it is adapted with $W_0=0$, if $W_t-W_s\sim N(0,t-s)$ for $0\le s\le t$, if $W_t-W_s$ is independent of $\mathcal F_s$, and if $t\mapsto W_t(\omega)$ is continuous.
Independence and the Gaussian law of increments fix all finite-dimensional distributions, and continuity of the paths is what separates Brownian motion from other processes with the same increments.
The scaled symmetric random walk replaces the Gaussian increments by coin flips: with $\xi_i=\pm1$ equally likely,
\begin{equation}\label{eq:scaledwalk}
 W^{(N)}_{t_k}=\sqrt{T/N}\sum_{i=1}^{k}\xi_i,\qquad t_k=kT/N,
\end{equation}
so that $\mathbb E\,W^{(N)}_{t_k}=0$ and $\operatorname{Var}W^{(N)}_{t_k}=t_k$ for every $N$.
Mean, variance and independence of increments therefore agree with Brownian motion for every $N$, whereas the lattice support and the non-Gaussian increments do not.
Donsker's theorem states that the whole path $W^{(N)}$ converges in law to $W$ as $N\to\infty$, which is the statement we test numerically.
This matters for pricing because the CRR tree of Section~\ref{sec:B} is exactly such a walk with a drift, so its convergence to geometric Brownian motion is a consequence of the convergence studied here.
